\section{Introduction}
\label{sec:intro}

Many practical uses of large language models embed the model in a small,
discrete environment that is best described as a text grid: a screenshot
serialised as ASCII, a board game state, an ARC-style puzzle, a tabular
form. In all of these settings the model must \emph{address} the grid
verbally---report the token at a cell, locate a feature, count
occurrences, identify the densest row. Whether this kind of addressing
works \emph{without} explicit step-by-step reasoning has direct
consequences for inference cost, latency, and the design of agentic
systems that ground decisions in serialised state.

A growing body of work argues that chain-of-thought (CoT) prompting often
hurts on in-context pattern-induction tasks~\citep{zheng2025curse,
sprague2025tocot}, and that text-only LLMs can pick up 2D-geometric
structure from a few in-context demonstrations~\citep{liu2025sparks}. A
parallel line of work shows that vision-language models fail on novel
polygons unless prompts label the geometry
explicitly~\citep{rudman2025polygons,wang2025vdlm}. And surface-level
repetition in the prompt is known to push next-token predictions toward
self-reinforced patterns~\citep{yan2024repetitions}. {\bf What none of
these threads has done is hold the geometry fixed, vary the amount of
in-context repetition about that same geometry, and ask whether direct
prompting can substitute for reasoning on verbal addressing of a
\emph{novel} shape.} That intersection is the gap we close.

\para{Our approach.} We design \addrtask, a controlled
$2{\times}2{\times}5$ factorial study (\figref{fig:rep_curves}). We
procedurally generate 80 novel 7$\times$7 ASCII grids, each containing
one irregular shape grown by a random walk over an alphabet of four
non-empty tokens. For every grid we issue one held-out
target addressing question drawn from five mechanically-checkable
question types (\qpoint, \qlookup, \qcount, \qrowmax, \qneighbor), padded
by $K \in \{0,1,2,4,8\}$ same-grid demonstration $(Q,A)$ pairs. We sweep
two models (\gpt, \haiku) and two prompting regimes (\Direct, \COT) at
each $K$, for 1{,}600 total queries. The design lets us cleanly factor
out the effect of reasoning, repetition, and query type with
within-item paired statistics.

\para{Headline result.} \Direct prompting plateaus at 79--93\% accuracy
even at $K{=}8$, while \COT lifts both models to 99--100\% from $K{=}0$
onward. A per-query-type breakdown reveals the cause: random-access
queries (\qpoint, \qlookup, \qneighbor) succeed at $\ge 95\%$ accuracy
under \Direct, while \qcount and \qrowmax---queries that require an
iterative scan of the grid---fail 12--68\% of the time. Paired McNemar
tests show the \Direct--\COT gap is significant at every
$(\text{model}, K)$ cell ($p \le 0.031$), and nine of ten survive a
Bonferroni correction over the ten paired tests. Repeating the same grid
helps \Direct prompting (\haiku: $0.65 \to 0.93$, $p<10^{-6}$) but
\emph{cannot} close the aggregation gap that \COT closes for free.

\para{Why this matters.} Our results refine the ``curse of CoT'' story
for in-context learning: the curse holds for rule induction but inverts
for verbal addressing of novel geometries, with the inversion driven by
exactly the scan-and-accumulate queries that \citet{sprague2025tocot}
identify as the regime where CoT actually buys accuracy. For
practitioners building LLM agents on top of grid-shaped state, enabling
CoT is a strictly cheaper reliability lever than padding the context
with same-grid demonstrations.

\para{Contributions.}
\begin{itemize}[leftmargin=*,itemsep=0pt,topsep=0pt]
    \item We propose \addrtask, a procedurally-generated benchmark of
    verbal addressing on novel 2D ASCII shapes that factors reasoning
    against in-context repetition with within-item paired statistics.
    \item We show that two production LLMs reach 99--100\% accuracy
    under \COT from $K{=}0$, while \Direct prompting plateaus at
    79--93\% even at $K{=}8$, and identify scan-based aggregation
    queries (\qcount, \qrowmax) as the sole source of the gap.
    \item We complement the result with a \miniarc sanity check showing
    \gpt exhibits the opposite \Direct${>}$\COT ordering on rule
    induction, reconciling our finding with the curse-of-CoT
    literature: the gap depends on whether the query is
    addressing or induction.
    \item We release code, prompts, raw model outputs, and analysis
    scripts at total reproducibility cost of about \$1.50.
\end{itemize}

\Secref{sec:related} surveys CoT skepticism, spatial reasoning in LLMs,
and in-context repetition effects. \Secref{sec:method} describes
\addrtask and the statistical analysis plan. \Secref{sec:results}
presents the main results. \Secref{sec:discussion} interprets them and
discusses limitations. \Secref{sec:conclusion} concludes.
